% TREŚĆ ROZDZIAŁU: zastąp wskazówki opisem testów i uzyskanych wyników.
% Każdą tabelę i każdy rysunek omów w tekście. Podawaj jednostki, warunki
% pomiaru oraz odwołanie do wymagania, które dany test weryfikuje.
\chapter{Weryfikacja i wyniki}
\label{chap:wyniki}

\section{Metodyka testów}
Opisz stanowisko, dane, warunki, procedurę i miary. Informacje powinny wystarczyć do powtórzenia najważniejszych testów.

\section{Wyniki}
Dane porównawcze obu algorytmów zawiera \cref{tab:porownanie-algorytmow}.
Zależność czasu sortowania bąbelkowego od liczby elementów pokazuje \cref{fig:wykres-sortowania}.
Kod funkcji wyszukującej maksimum zamieszczono w \cref{lst:max-element}.
\subsection{Porównanie algorytmów}
Zestawienie czasu i liczby porównań dla dwóch algorytmów zawiera
\cref{tab:porownanie-algorytmow}. Są to dane dydaktyczne, a nie wyniki pomiarów.

\begin{table}[htbp]
  \centering
  \caption{Porównanie przykładowych wyników dwóch algorytmów sortowania}
  \label{tab:porownanie-algorytmow}
  \begin{tabularx}{\textwidth}{@{}c *{4}{>{\centering\arraybackslash}X}@{}}
    \toprule
    Liczba elementów $n$
      & Bubble Sort czas [ms]
      & Insertion Sort czas [ms]
      & Bubble Sort porównania
      & Insertion Sort porównania \\
    \midrule
    100  & \num{0.4}   & \num{0.2}   & \num{4950}     & \num{2600}    \\
    500  & \num{8.9}   & \num{4.1}   & \num{124750}   & \num{65000}   \\
    1000 & \num{35.2}  & \num{16.3}  & \num{499500}   & \num{260000}  \\
    2000 & \num{141.5} & \num{65.1}  & \num{1999000}  & \num{1040000} \\
    5000 & \num{883.7} & \num{406.4} & \num{12497500} & \num{6500000} \\
    \bottomrule
  \end{tabularx}
\end{table}

Zależność czasu sortowania bąbelkowego od liczby elementów pokazuje
\cref{fig:wykres-sortowania}. Wykres dotyczy wyłącznie Bubble Sort.

\begin{figure}[htbp]
  \centering
  \includegraphics[width=0.8\textwidth]{%
    chapters/05-weryfikacja-i-wyniki/figures/wykres_sortowania.png}
  \caption{Przykładowa zależność czasu sortowania od liczby elementów}
  \label{fig:wykres-sortowania}
\end{figure}

\section{Ocena spełnienia wymagań}
Dla każdego wymagania podaj wynik testu oraz status: spełnione, częściowo spełnione albo niespełnione. Odnieś się do kryteriów z \cref{chap:cel-wymagania}.

\section{Dyskusja}
Zinterpretuj rezultaty, wyjaśnij odstępstwa i wskaż ograniczenia wiarygodności oceny.

\section{Wykresy tworzone w LaTeX}
\label{sec:wykresy-latex}

\subsection{Wykres funkcji matematycznej}

Na przedziale \([-2, 4]\) funkcja \(f(x) = x^{2} - 2x - 3\) jest parabolą
skierowaną ramionami do góry. Pochodna \(f'(x) = 2x - 2\) zeruje się
dla \(x = 1\), więc minimum wynosi \(f(1) = -4\). Ten punkt jest
zaznaczony na wykresie.

\begin{center}
  \begin{tikzpicture}
    \begin{axis}[
        width=0.82\textwidth,
        height=7.2cm,
        xlabel={$x$},
        ylabel={$f(x)$},
        grid=both,
        minor tick num=1,
        tick label style={font=\small},
        label style={font=\small},
        scaled ticks=false,
        ymin=-5, ymax=6,
        enlarge y limits=false,
        clip=false,
      ]
      \addplot[domain=-2:4, samples=101, very thick, projectblue]
        {x^2 - 2*x - 3};
      \addplot[only marks, mark=*, mark size=2.2pt, red!70!black]
        coordinates {(1,-4)};
      \node[anchor=west, font=\small, inner sep=1pt]
        at (axis cs:1.15,-4.55) {$(1,-4)$};
    \end{axis}
  \end{tikzpicture}
\end{center}

\subsection{Porównanie czasów wykonania}

\begin{figure}[htbp]
  \centering
  \begin{tikzpicture}
    \begin{axis}[
        width=0.88\textwidth,
        height=7.6cm,
        xlabel={Liczba elementów $n$},
        ylabel={Czas wykonania [ms]},
        grid=major,
        legend pos=north west,
        legend cell align={left},
        tick label style={font=\small},
        label style={font=\small},
        legend style={font=\small},
        scaled ticks=false,
        xmin=0, xmax=5400,
        ymin=0, ymax=1000,
      ]
      \addplot[
          very thick,
          projectblue,
          mark=*,
          mark size=1.8pt,
        ] table[col sep=comma, header=true, x=n, y=bubble_ms]
          {dane_porownawcze.csv};
      \addplot[
          thick,
          dashed,
          orange!75!black,
          mark=square*,
          mark size=2.1pt,
        ] table[col sep=comma, header=true, x=n, y=insertion_ms]
          {dane_porownawcze.csv};
      \legend{Bubble Sort, Insertion Sort}
    \end{axis}
  \end{tikzpicture}
  \caption{Przykładowe czasy wykonania algorytmów sortowania}
  \label{fig:czasy-sortowania}
\end{figure}

\Cref{fig:czasy-sortowania} pokazuje, że obie serie rosną szybciej niż liniowo,
a czas Bubble Sort jest w każdym punkcie około dwukrotnie większy niż
czas Insertion Sort. Wraz ze wzrostem liczby elementów różnica się
powiększa: dla \(n = 5000\) czasy wynoszą odpowiednio
\qty{883.7}{\milli\second} oraz \qty{406.4}{\milli\second}.
Podane wartości są syntetyczne i nie stanowią wyników pomiarów.
